% Vietnamese translation for OI Public Library
% Source-File: IOI中国国家候选队论文/2005/杨俊/杨俊.ppt
\documentclass[11pt]{article}
\usepackage{oipl-vn}

\title{Bản trình chiếu: Chiến lược nhị phân}
\author{Yang Jun}
\date{2005}

\begin{document}
\maketitle

\origtitle{Slide companion for binary-search strategy}
\sourcepath{IOI China candidate team papers / 2005 / Yang Jun / Yang Jun.ppt}

\section{Phạm vi bản dịch}

Bản trình chiếu đi cùng bài viết về chiến lược nhị phân. Các slide đọc được
nhắc tới WinterCamp 2005, tập \(S=\{3,7,2,6,8,1,5\}\), đáp án \(5\), ví dụ
\(N=8,K=2\), và các độ phức tạp \(O(n)\), \(O(N\log^2 N)\), \(O(n\log n)\).

\section{Tư tưởng}

Chiến lược nhị phân không chỉ là tìm kiếm nhị phân trên mảng đã sắp. Nó là
cách chia không gian lời giải thành hai phần bằng một điều kiện đơn điệu:
nếu một giá trị thỏa, mọi giá trị về một phía cũng thỏa; nếu không, ta bỏ
phía đó.

\section{Ứng dụng}

Các dãy bit \(0,0,\ldots,1,1,\ldots\) trong slide minh họa hàm kiểm tra đơn
điệu. Khi bài toán tối ưu có thể chuyển thành câu hỏi ``có tồn tại lời giải
với tham số \(x\) hay không'', ta nhị phân trên \(x\), còn phần kiểm tra dùng
thuật toán riêng.

\section{Kết luận}

Bài trình bày nhấn mạnh hai bước: chứng minh tính đơn điệu, rồi thiết kế hàm
kiểm tra đủ nhanh. Nếu thiếu một trong hai bước, tìm kiếm nhị phân sẽ không
đem lại lời giải đúng.

\end{document}
